\documentclass[10pt]{article}
\usepackage[margin=.8in]{geometry}
\usepackage{amsmath,hyperref}
\title{Guide: Mersenne prefix exits}
\author{Collatz Lean research project}
\date{September 6, 2026}
\begin{document}
\maketitle
\section*{What changed}
A growing prefix previously prevented a fixed search depth from handling all Mersenne parameters. We now express the states after that prefix as
\[
 X-7,\quad 3X-10,\qquad X=3^{4h+5},\quad k=6h+4.
\]
This permits searching the exponent modulo a power of two, instead of simulating a separately bounded prefix for each k.
\section*{The first formal rule}
For every $k=1468+1536z$, the reduced inverse pair merges after $k+12$ steps. The additional twelve steps are the same symbolic rule throughout the progression. The prefix length is unbounded.

This is useful for the transfer program: a smaller transfer instance at
$v=8w+5$, where $w=(2^k-7)/9$, can prove the instance at $u=2^k-1$.
It is a conditional step, not a proof that every Collatz orbit reaches one.
\section*{Where to inspect it}
\raggedright
The formal source is \texttt{lean/Collatz/MersenneExit.lean}.
The declarations \texttt{mersenne\_exit\_merge},
\texttt{mersenne\_inverse\_exit}, and
\texttt{mersenne\_inverse\_exit\_progression} separate the small merge,
the growing prefix, and their composition.
The declarations \texttt{mersenne\_exit\_smaller} and
\texttt{AffineTransfer.mersenne\_exit} verify strict decrease and the conditional induction step.

The experiment and exact saved census are
\texttt{analysis/mersenne\_exit\_search.py} and
\texttt{analysis/mersenne\_exit\_results.json}.
The scope and derivation are recorded in
\texttt{analysis/MERSENNE\_EXIT\_CHECK.md}.
\section*{Reproduce}
From the repository root, run
\begin{verbatim}
python3 analysis/mersenne_exit_search.py --depth 16
python3 -m unittest discover -s analysis -p 'test_mersenne*.py'
python3 analysis/audit_theorems.py --build
\end{verbatim}
From \texttt{lean/}, run \texttt{lake build Collatz.MersenneExit}.
Lean and Lake must use the repository's pinned toolchain.
\section*{What remains}
The depth-16 census covers 110 of 4096 eligible exponent residues. The unresolved classes are not counterexamples. A full proof still needs transfers or reductions that cover every remaining case and make the induction well founded. The finite census has not been imported as a Lean theorem.
\end{document}
